\documentclass[11pt,a4paper]{article}

\input{glyphtounicode}
\pdfgentounicode=1

\usepackage[T1]{fontenc}
\usepackage[utf8]{inputenc}
\usepackage{lmodern}
\usepackage[a4paper,margin=2.4cm]{geometry}
\usepackage{microtype}
\usepackage{xcolor}
\usepackage{hyperref}
\usepackage{longtable}
\usepackage{booktabs}
\usepackage{array}
\usepackage{listings}
\IfFileExists{upquote.sty}{\usepackage{upquote}}{}

\hypersetup{
  colorlinks=true,
  linkcolor=blue!50!black,
  urlcolor=blue!50!black,
  pdftitle={Guia de manutencao do site MFQ},
  pdfauthor={Codex}
}

\setlength{\parindent}{0pt}
\setlength{\parskip}{6pt}
\renewcommand{\contentsname}{Sumário}

\definecolor{codebg}{RGB}{247,248,250}
\definecolor{coderule}{RGB}{210,214,220}

\lstdefinestyle{mfqstyle}{
  basicstyle=\ttfamily\small,
  backgroundcolor=\color{codebg},
  frame=single,
  rulecolor=\color{coderule},
  framerule=0.5pt,
  columns=fullflexible,
  keepspaces=true,
  showstringspaces=false,
  breaklines=false,
  tabsize=2,
  aboveskip=0.8em,
  belowskip=0.8em
}

\lstnewenvironment{terminal}{
  \lstset{style=mfqstyle}
}{}

\newcommand{\repo}{\path{mathfoundq}}
\newcommand{\mfqdir}{\path{mfq}}

\title{Guia Completo de Manutenção do Site MFQ}
\author{}
\date{25 de março de 2026}

\begin{document}

\maketitle

\section*{Cheat Sheet}

Todos os comandos deste documento aparecem sem \texttt{\$} no começo, para facilitar copiar e colar.
Os comandos abaixo assumem que você já entrou na pasta raiz do repositório.

\subsection*{Fluxo diário}

\begin{terminal}
git status
git pull --rebase
cd mfq
hugo server
\end{terminal}

\subsection*{Gerar o site final}

\begin{terminal}
cd mfq
hugo
\end{terminal}

O site gerado ficará em \path{mfq/public/}.

\subsection*{Publicar no servidor}

Confirme host, usuário e credenciais com a equipe antes de usar.

\begin{terminal}
cd mfq
hugo
cd public
scp -r * mfq@ssh.ime.unicamp.br:public_html/
\end{terminal}

\subsection*{Arquivos que você mais vai editar}

\begin{longtable}{>{\raggedright\arraybackslash}p{0.28\textwidth} p{0.66\textwidth}}
\toprule
Arquivo ou pasta & Uso \\
\midrule
\endhead
\path{mfq/content/english/people/} & perfis das pessoas \\
\path{mfq/content/english/publications/publications.yaml} & lista de publicações \\
\path{mfq/content/english/talks/upcoming.md} & talks futuras \\
\path{mfq/content/english/talks/previous.md} & talks passadas \\
\path{mfq/content/english/research/} & tópicos de pesquisa \\
\path{mfq/content/english/gallery/index.md} & eventos e fotos da galeria \\
\path{mfq/data/en/homepage.yml} & cards de pesquisa e textos da home \\
\bottomrule
\end{longtable}

\tableofcontents
\newpage

\section{Visão geral do repositório}

O site é um site estático em Hugo. O trabalho do dia a dia acontece quase sempre nestes lugares:

\begin{itemize}
  \item \path{mfq/content/english/people/}: uma pasta por pessoa, com \path{index.md} e a foto.
  \item \path{mfq/content/english/publications/}: página de publicações e arquivo \path{publications.yaml}.
  \item \path{mfq/content/english/research/}: uma página por tema de pesquisa.
  \item \path{mfq/content/english/talks/}: introdução, upcoming e previous.
  \item \path{mfq/content/english/gallery/}: galeria e eventos com imagens.
  \item \path{mfq/data/en/homepage.yml}: textos e cards da página inicial.
\end{itemize}

O tema do Hugo já está dentro do próprio projeto, em \path{mfq/themes/vex-hugo}. Para rodar o site localmente,
use \texttt{hugo server} dentro de \path{mfq/}.

\section{Instalação}

\subsection{Windows}

Se você está no Windows, o procedimento recomendado para manter este site é o seguinte:

\begin{enumerate}
  \item Faça backup dos seus arquivos.
  \item Instale Linux.
  \item Remova o Windows.
  \item Volte para a próxima subseção, \textbf{Linux}, e siga normalmente.
\end{enumerate}

\subsection{Linux}

\subsubsection{Git}

Escolha apenas o bloco da sua distribuição Linux.

\textbf{Debian / Ubuntu}
\begin{terminal}
sudo apt update
sudo apt install git
\end{terminal}

\textbf{Fedora}
\begin{terminal}
sudo dnf install git
\end{terminal}

\textbf{Arch / Manjaro}
\begin{terminal}
sudo pacman -S git
\end{terminal}

Depois configure seu nome e seu e-mail uma vez:

\begin{terminal}
git config --global user.name "Seu Nome"
git config --global user.email "seu.email@exemplo.com"
\end{terminal}

\subsubsection{Hugo}

Este repositório compila com Hugo Extended. Instale a versão ``extended''.

\textbf{Debian / Ubuntu}
\begin{terminal}
sudo apt update
sudo apt install hugo
\end{terminal}

\textbf{Fedora}
\begin{terminal}
sudo dnf install hugo
\end{terminal}

\textbf{Arch / Manjaro}
\begin{terminal}
sudo pacman -S hugo
\end{terminal}

Confira se os programas foram instalados:

\begin{terminal}
git --version
hugo version
\end{terminal}

\section{Primeira configuração do projeto}

\subsection{Se você for o mantenedor principal}

A primeira coisa a fazer é criar um \emph{fork} do repositório no GitHub.
Esse fork é uma cópia do projeto na sua conta, o que facilita continuidade, administração e recuperação do histórico.

Depois, clone o seu fork e entre na pasta:

\begin{terminal}
git clone <URL_DO_SEU_FORK> mathfoundq
cd mathfoundq
\end{terminal}

\subsection{Se você só vai fazer edições pontuais}

Se você não for o mantenedor principal, peça à equipe a URL do repositório que deve receber as mudanças e clone normalmente:

\begin{terminal}
git clone <URL_DO_REPOSITORIO> mathfoundq
cd mathfoundq
\end{terminal}

\noindent O branch principal deste projeto é \texttt{master}.

\section{Como usar Git neste projeto}

\subsection{O que é o Git}

Git é o sistema que registra o histórico do projeto.
Ele permite saber o que mudou, quando mudou e quem mudou.
Na prática, é ele que organiza as edições do site sem transformar tudo em uma pasta caótica de cópias com nomes como
\texttt{site-final-agora-v2-definitivo}.

\subsection{Conceitos que importam aqui}

\begin{itemize}
  \item \textbf{Repositório remoto}: a versão do projeto que está no GitHub.
  \item \textbf{Fork}: uma cópia desse repositório remoto na sua própria conta do GitHub.
  \item \textbf{Repositório local}: a cópia que fica no seu computador, dentro da pasta \repo.
  \item \textbf{Arquivo modificado}: um arquivo que você editou localmente e que o Git percebeu como diferente da última versão salva.
  \item \textbf{Commit}: um ``salvamento com histórico'', isto é, um ponto registrado no tempo com uma mensagem explicando a mudança.
\end{itemize}

Quando você altera um arquivo como \path{publications.yaml} ou \path{index.md}, a mudança existe primeiro só no seu
repositório local. Ela só entra no histórico do Git quando você faz um commit. Ela só vai para o GitHub quando você faz um push.

\subsection{Fluxo simples no terminal}

Você não precisa rodar todos os comandos possíveis do Git toda vez.
Na maior parte do tempo, um fluxo simples resolve:

\begin{enumerate}
  \item Entre na pasta raiz do repositório.
  \item Rode \texttt{git pull --rebase} para atualizar sua cópia local.
  \item Faça suas edições.
  \item Rode \texttt{git status} para ver o que mudou.
  \item Faça um commit.
  \item Envie com \texttt{git push}.
\end{enumerate}

Um exemplo mínimo:

\begin{terminal}
git pull --rebase
git status
git add .
git commit -m "Atualiza site MFQ"
git push
\end{terminal}

Se quiser revisar linha por linha antes de salvar, aí sim vale usar \texttt{git diff}, mas isso é opcional no dia a dia.

\subsection{Fluxo simples no VS Code}

Se você usa VS Code, dá para fazer quase tudo por interface gráfica.

\begin{enumerate}
  \item Abra a pasta do projeto no VS Code.
  \item Vá no painel \textbf{Source Control}.
  \item Veja a lista de arquivos modificados.
  \item Clique em um arquivo para ver o diff visual.
  \item Use o botão \texttt{+} para incluir a mudança no commit.
  \item Escreva a mensagem do commit.
  \item Clique em \textbf{Commit}.
  \item Use os botões de \textbf{Sync}, \textbf{Pull} e \textbf{Push} quando precisar.
\end{enumerate}

O suporte básico a Git já vem no próprio VS Code. Extensões de Git podem ajudar, mas você não precisa delas para o essencial.

\subsection{Comandos que vale conhecer}

\begin{longtable}{>{\raggedright\arraybackslash}p{0.26\textwidth} p{0.68\textwidth}}
\toprule
Comando & O que faz \\
\midrule
\endhead
\texttt{git pull --rebase} & atualiza sua cópia local com base no remoto configurado \\
\texttt{git status} & mostra arquivos alterados \\
\texttt{git diff} & mostra exatamente o que mudou; é útil, mas não obrigatório toda vez \\
\texttt{git add <arquivo>} & coloca um arquivo no próximo commit \\
\texttt{git add .} & coloca todas as alterações do diretório atual \\
\texttt{git commit -m "..."} & cria um commit com mensagem \\
\texttt{git push} & envia seus commits para o remoto configurado \\
\bottomrule
\end{longtable}

\section{Como usar Hugo neste projeto}

Entre em \path{mfq/} para rodar os comandos do Hugo:

\begin{terminal}
cd mfq
\end{terminal}

\subsection{Visualizar o site localmente}

\begin{terminal}
hugo server
\end{terminal}

Depois abra no navegador a URL local mostrada pelo próprio Hugo no terminal.

\subsection{Gerar a versão final do site}

\begin{terminal}
hugo
\end{terminal}

Isso gera a pasta \path{mfq/public/}, que contém o site estático pronto para ser publicado.

\subsection{O que verificar antes de publicar}

\begin{itemize}
  \item Se a pessoa nova aparece na seção correta.
  \item Se a imagem carregou.
  \item Se a publicação entrou no ano correto.
  \item Se links de autores e links externos funcionam.
  \item Se a talk foi para \emph{upcoming} ou \emph{previous} corretamente.
  \item Se o tema de pesquisa aparece em \path{/research}.
\end{itemize}

\section{Adicionar uma pessoa nova}

Cada pessoa é um \emph{page bundle} do Hugo: uma pasta própria com \path{index.md} e a imagem dentro dela.

\subsection{Passo a passo}

\begin{enumerate}
  \item Crie uma pasta em \path{mfq/content/english/people/}.
  \item O nome da pasta será o \emph{slug} da pessoa. Exemplo: \path{maria_silva/}.
  \item Coloque a foto dentro da mesma pasta.
  \item Crie um arquivo \path{index.md}.
  \item Escolha a categoria correta: \texttt{leader}, \texttt{researcher}, \texttt{student} ou \texttt{alumni}.
  \item Deixe \texttt{draft: false}, senão a página não aparece.
\end{enumerate}

\subsection{Estrutura esperada}

\begin{terminal}
mfq/content/english/people/maria_silva/
mfq/content/english/people/maria_silva/index.md
mfq/content/english/people/maria_silva/maria_silva.jpg
\end{terminal}

\subsection{Modelo de index.md}

\begin{terminal}
---
title: "Maria Silva"
images:
  - "maria_silva.jpg"
summary: Ph.D. student
categories: student
draft: false
---

Short biography in English.

[Lattes](http://exemplo.com)
[ORCID](https://orcid.org/0000-0000-0000-0000)
\end{terminal}

\subsection{Observações importantes}

\begin{itemize}
  \item A foto pode ficar na resolução original. O Hugo corta e gera as versões usadas no site.
  \item O \texttt{summary} é o texto curto que aparece nos cards.
  \item O corpo do arquivo, abaixo do front matter, é o texto da página da pessoa.
  \item O nome da pasta é importante porque ele pode ser usado em publicações como link do autor.
\end{itemize}

\section{Adicionar uma publicação nova}

As publicações não ficam em um arquivo Markdown. Elas ficam em \path{mfq/content/english/publications/publications.yaml}.

\subsection{Campos principais}

Cada item precisa ter:

\begin{itemize}
  \item \texttt{title}
  \item \texttt{link}
  \item \texttt{authors}
  \item \texttt{published\_in}
  \item \texttt{date}
\end{itemize}

\subsection{Modelo de item}

\begin{terminal}
- title: Example paper title
  link: https://doi.org/10.0000/exemplo
  authors:
  - name: Maria Silva
    page: maria_silva
  - name: Coauthor Name
  published_in: Journal Name
  date: '2026-03-25'
\end{terminal}

\subsection{Regras práticas}

\begin{itemize}
  \item Use a data no formato \texttt{AAAA-MM-DD}.
  \item A página de publicações é ordenada automaticamente por data decrescente.
  \item Se você colocar \texttt{page: maria\_silva}, o nome do autor vira link para \path{/people/maria_silva/}, desde que essa pasta exista.
  \item Se o autor não tiver página no site, use só \texttt{name}.
\end{itemize}

\section{Adicionar um evento novo na galeria}

Hoje o site não tem uma coleção separada chamada ``events''. Os eventos aparecem na galeria, em
\path{mfq/content/english/gallery/index.md}, com as imagens na mesma pasta.

\subsection{Passo a passo}

\begin{enumerate}
  \item Copie a imagem para \path{mfq/content/english/gallery/}.
  \item Edite \path{mfq/content/english/gallery/index.md}.
  \item Adicione um bloco novo com título, imagem, data, local e descrição.
\end{enumerate}

\subsection{Exemplo}

\begin{terminal}
## 2026

### **Workshop on Quantum Foundations**

**Date:** March 20 to March 22, 2026
**Location:** Campinas, SP, Brazil
**Description:** Short description of the event.

<hr>
\end{terminal}

\subsection{Observações}

\begin{itemize}
  \item A imagem precisa estar na mesma pasta de \path{index.md}.
  \item O shortcode \texttt{\{\{< image "..."\ >\}\}} é o jeito correto de inserir a imagem.
  \item Se quiser manter a galeria organizada, coloque o evento dentro do ano correspondente.
\end{itemize}

\section{Alterar ou adicionar tópicos de pesquisa}

\subsection{Editar um tópico existente}

Cada tópico é um arquivo em \path{mfq/content/english/research/}. Exemplos atuais:

\begin{itemize}
  \item \path{mfq/content/english/research/nonlocality.md}
  \item \path{mfq/content/english/research/contextuality.md}
  \item \path{mfq/content/english/research/gpts.md}
\end{itemize}

Edite o texto, o resumo e a imagem do arquivo desejado.

\subsection{Adicionar um tópico novo}

Crie um arquivo novo em \path{mfq/content/english/research/}. Exemplo:

\begin{terminal}
---
title: "Quantum Steering"
date: 2026-01-01T11:22:00+00:00
image: "images/steering.png"
description: "detailed view of quantum steering research"
summary: Short summary shown on the research list page.
draft: false
---

Full page text in English.
\end{terminal}

\subsection{O que controlar}

\begin{itemize}
  \item \texttt{draft: false} para a página aparecer.
  \item O campo \texttt{date} ajuda a controlar a ordem de exibição dos tópicos.
  \item O campo \texttt{summary} é o texto resumido da página \path{/research}.
\end{itemize}

\subsection{Importante: home e página de research não são a mesma coisa}

Adicionar um arquivo novo em \path{mfq/content/english/research/} faz o tema aparecer na página
\path{/research}. Mas os cards da home são controlados separadamente em
\path{mfq/data/en/homepage.yml}.

Se quiser que o novo tema também apareça na home, adicione um item novo em \path{about_product.items}:

\begin{terminal}
- image: "images/steering.png"
  title: "Quantum Steering"
  content: "Short summary for the homepage."
  button:
    enable: true
    label: "Read more"
    link: "research/steering"
\end{terminal}

\section{Adicionar uma Nonlocal Talk}

A página \path{/talks} hoje é montada juntando os arquivos da pasta \path{mfq/content/english/talks/}.
No uso normal, você deve editar estes arquivos:

\begin{itemize}
  \item \path{introduction.md}
  \item \path{upcoming.md}
  \item \path{previous.md}
\end{itemize}

\subsection{Fluxo recomendado}

\begin{enumerate}
  \item Coloque a próxima talk em \path{upcoming.md}.
  \item Depois que a talk acontecer, mova o bloco para o topo de \path{previous.md}.
  \item Se não houver talks agendadas, deixe \path{upcoming.md} com a mensagem ``More Soon!''.
\end{enumerate}

\subsection{Modelo para uma talk futura}

\begin{terminal}
<h1 style="display:inline-block;">Speaker Name</h1>
&nbsp;&nbsp;&nbsp;
<h2 style="display:inline-block;">(Institution, Country)</h2>

## Talk Title
### Apr 10, 2026

Abstract of the talk.

<br><br><hr>
\end{terminal}

\subsection{Observações}

\begin{itemize}
  \item O conteúdo mistura Markdown com HTML. Mantenha esse padrão.
  \item Não crie um arquivo novo por talk, a menos que você também vá alterar o template.
  \item A ordem da página depende dos arquivos existentes na seção, então o fluxo mais simples é continuar usando \path{upcoming.md} e \path{previous.md}.
\end{itemize}

\section{Ajustes frequentes fora desses blocos}

\subsection{Textos gerais da home}

Arquivo principal: \path{mfq/data/en/homepage.yml}

Lá ficam, por exemplo:

\begin{itemize}
  \item título e texto do banner inicial;
  \item seção ``A quick glance at our research'';
  \item links dos botões da home.
\end{itemize}

\subsection{Página About / Contact}

Arquivo principal: \path{mfq/content/english/about/_index.md}

Ali você ajusta:

\begin{itemize}
  \item texto institucional;
  \item e-mail do grupo;
  \item localização;
  \item conteúdo da página de contato.
\end{itemize}

\section{Checklist antes de fechar a tarefa}

\begin{itemize}
  \item Rode \texttt{git diff} e confira se você editou só o que queria.
  \item Rode \texttt{hugo} dentro de \path{mfq/}.
  \item Revise links, imagens e ortografia.
  \item Faça commit com mensagem clara.
\end{itemize}

\section{Publicação no servidor}

O deploy é feito copiando o conteúdo de \path{mfq/public/} para
\path{public_html} no servidor do IMECC. Como esse passo depende de credenciais e política atual
de acesso, confirme com a equipe antes de publicar.

\begin{terminal}
cd mfq
hugo
cd public
scp -r * mfq@ssh.ime.unicamp.br:public_html/
\end{terminal}

Se esse procedimento mudar no futuro, atualize primeiro esta seção do manual.

\section{Problemas comuns}

\subsection{A página não apareceu}

Verifique:

\begin{itemize}
  \item se \texttt{draft: false} está no front matter;
  \item se o arquivo foi criado na pasta correta;
  \item se o nome da imagem bate com o nome usado no arquivo;
  \item se o \texttt{hugo} terminou sem erro.
\end{itemize}

\subsection{O link do autor na publicação não funciona}

Verifique se o valor de \texttt{page} no YAML é exatamente o nome da pasta da pessoa em
\path{mfq/content/english/people/}.

\subsection{O preview local não abriu}

Use:

\begin{terminal}
cd mfq
hugo server
\end{terminal}

Se ainda falhar, verifique se você está dentro de \path{mfq/} e se o Hugo está instalado corretamente com
\texttt{hugo version}.

\section{Resumo final}

Se você lembrar só de quatro coisas, lembre destas:

\begin{enumerate}
  \item Entre em \path{mfq/} para rodar o Hugo.
  \item Pessoas ficam em \path{people/}, publicações em \path{publications.yaml}, talks em \path{talks/}.
  \item Para a home, muitas vezes você também precisa mexer em \path{mfq/data/en/homepage.yml}.
  \item Sempre gere o site com \texttt{hugo} antes de publicar.
\end{enumerate}

\end{document}
